\subsection{Three further full two-torsion exclusions}
\begin{proposition}\label{prop:decic238-240-250}
Let $K$ be a number field of degree $10$. No elliptic curve over $K$
contains any of the groups
\[
 C_2\times C_{38},\qquad C_2\times C_{40},\qquad C_2\times C_{50}.
\]
Consequently, no elliptic curve over $K$ has a point of order
$76$, $80$, or $100$.
\end{proposition}
\begin{proof}
For $n=38,40,50$, put $X_n=X_1(2,n)$. A point over $K$ gives a
noncuspidal effective rational divisor $D$ of degree $10$ on $X_n$.
We allow repeated support, so this also covers points defined over
proper subfields of $K$.

The genera of these curves are $64,57,109$, respectively. Their
Jacobians have rank zero. Our certificates verify this directly:
for every even character trivial on the corresponding subgroup,
the winding vector generates the full cuspidal modular-symbol space
under Hecke operators. The dimensions, weighted by the character-field
degrees, sum to the stated genus. The indices of the congruence
subgroups in $\operatorname{PSL}_2(\mathbb Z)$ are $1080,1152,1800$.
Thus Abramovich's bound with the Kim--Sarnak constant gives
\[
 \operatorname{gon}_{\overline{\Q}}X_n
 \geq\frac{325}{32768}[\operatorname{PSL}_2(\mathbb Z):\overline\Gamma_n]
 >10.
\]
All three curves have good reduction at $3$. Their cusps split over
cyclotomic extensions unramified at $3$. By
Lemma~\ref{lem:decic-cuspidal-reduction}, the reduction of $D$ cannot
be entirely cuspidal.

We describe the necessary divisor-class test. Let $f:X_n\to Y$ be
one of the maps used below, choose a rational cusp $c_0$ on $Y$, and put
\[
 H_Y=\operatorname{red}_3 J_Y(\Q)\subseteq J_Y(\mathbb F_3).
\]
Write $\overline D=N+C$, where $N$ has noncuspidal support and $C$
has cuspidal support. A necessary condition is
\begin{equation}\label{eq:decic-small-picard}
 [f_*N-(\deg N)c_0]\in
 \bigcup_{C'}\bigl(H_Y-[C'-(\deg C')c_0]\bigr),
\end{equation}
where $C'$ runs over all effective cuspidal divisors on $Y_{\mathbb F_3}$
of degree $10-\deg N$. Allowing all such divisors enlarges the possible
images of $C$ and is therefore valid for exclusion. If a source point
has degree $d$ and its image has degree $e$, its pushforward has
multiplicity $d/e$. Every test includes this multiplicity.

The rational Jacobian groups used in this test are generated by
rational cuspidal divisors. For the cyclic targets this follows from
\cite[Theorem~4.13 and Corollary~4.14]{DEHMZ}; for $X_1(2,20)$ it
follows from \cite[Proposition~4.16 and Table~3]{DEHMZ}.
The relevant finite groups and rational images are
\[
\begin{array}{c|c|c|r}
 Y&J_Y(\mathbb F_3)&H_Y&[J_Y(\mathbb F_3):H_Y]\\\hline
 X_1(19)&[8766]&[4383]&2\\
 X_1(2,20)&[8,120,120]&[4,60,120]&4\\
 X_1(25)&[2503105]&[227555]&11\\
 X_1(40)&[2,2,2,2,2,16,68880,895440]
          &[2,2,2,8,120,895440]&4592
\end{array}
\]
Here brackets denote invariant factors. The computations determine the
embedded subgroups, not only their abstract structures.

We first enumerate every noncuspidal closed point of degree at most
$10$ over $\mathbb F_3$. The numbers of such points are
\[
\begin{array}{c|rrrrrrr}
 n\backslash d&4&5&6&7&8&9&10\\\hline
 38&0&54&63&162&756&2270&6129\\
 40&32&96&64&480&504&2048&6000\\
 50&0&0&0&60&1170&1860&5040
\end{array}
\]
There are none in degrees $1,2,3$. The enumeration starts with every
elliptic-curve isomorphism class over the finite fields and every
admissible torsion basis. Independent point counts from the Frobenius
polynomials verify completeness in each degree. We then take orbits
under all level automorphisms. Their orders are $54,32,60$,
respectively. The numbers of orbits in degrees $4,\ldots,10$ are
\[
\begin{array}{c|rrrrrrr}
38&0&1&2&3&14&43&116\\
40&1&3&2&15&21&64&195\\
50&0&0&0&1&21&31&84.
\end{array}
\]
For a mixed divisor, only its first support is normalized. We retain
every actual point for the second support, including repetitions.
The possible mixed types are $(5,5)$ for $n=38$, and
$(4,4),(4,5),(4,6),(5,5)$ for $n=40$. There are no mixed types for
$n=50$. This gives $233,781,137$ patterns, respectively. A pattern
specifies $N$; its cusp complement is treated by
\eqref{eq:decic-small-picard}.

\medskip\noindent
\emph{The case $n=38$.}
Put $V=2Q$, where $(P,Q)$ is the universal level structure on $X_{38}$.
We use four maps to $X_1(19)$: retain $V$, or quotient by one of
\[
 \langle P\rangle,\quad\langle19Q\rangle,
 \quad\langle P+19Q\rangle
\]
and retain the image of $V$. Applying
\eqref{eq:decic-small-picard} leaves $35$ patterns: one of degree $8$,
two of degree $9$, fourteen of degree $10$, and eighteen of type $(5,5)$.

Now use the forgetful map $\pi:X_{38}\to X_1(38)$, $(E,P,Q)\mapsto(E,Q)$.
For a prime $q$ of good reduction put
\[
 A_q=T_q-1-q\langle q\rangle.
\]
The Eichler--Shimura relation and torsion injectivity at the odd prime
$q$ show that $A_q$ annihilates $J_1(38)(\Q)$. The cusps of
$X_1(38)_{\mathbb F_3}$ have degrees $1^{18},9^2$. Every cusp in a
complement of degree at most $5$ is therefore rational over
$\mathbb F_3$. The Tate cusp calculation shows that each is already
rational over $\Q$. Consequently $A_5$ kills such cusp divisors as
divisors: on a rational cusp $c$ one has
$T_5c=c+5\langle5\rangle c$.
It follows that $A_5(\pi_*N)$ must be principal. Exact Riemann--Roch
tests show that it is nonprincipal for all $35$ patterns. A separate
calculation using the native Hecke operator verifies every test.
Thus this case is impossible.

\medskip\noindent
\emph{The case $n=40$.}
The maps to $X_1(2,20)$ are
\[
 (E,P,Q)\longmapsto(E,P,2Q)
\]
and
\[
 (E,P,Q)\longmapsto
 (E/\langle20Q\rangle,\phi(P),\phi(Q)).
\]
The first leaves $288$ of the $781$ patterns, and the second leaves
$118$. We then apply the forgetful map to $X_1(40)$ and test
\eqref{eq:decic-small-picard} in its full finite Picard group.
Its cusps over $\mathbb F_3$ have degrees $1^{12},2^6,4^6$.
All $118$ patterns fail, even after allowing every cuspidal complement.
The finite Picard presentation is proved with \texttt{Proof:=true}.
It is independent of the degree being studied; we copy it into the
present directory with its defining function field, factor-base
geometry and provenance. All new divisor images are checked by exact
principal ideal identities in both the finite and infinite maximal
orders. A separate reconstruction verifies all $146$ images used in
the eight disjoint parts of the sieve. Hence this case is impossible.

\medskip\noindent
\emph{The case $n=50$.}
We use the four maps to $X_1(25)$ obtained by retaining $2Q$, either
on $E$ or on its quotients by the three rational subgroups of order $2$.
The class-group test leaves three patterns, each consisting of one
closed point of degree $8$ and a cusp complement of degree $2$.
Under the forgetful map to $X_1(50)$, the cusp complement consists
of globally rational cusps: the finite cusp degrees are
$1^{20},4^4,10^2$, and the Tate calculation identifies all rational
residue cusps with rational global cusps.

The same argument as for $n=38$ shows that
$A_7(\pi_*N)$ must be principal. The three exact Riemann--Roch tests
are all negative. The positive Hecke divisor has degree $64$ in
each case. We independently verify the Hecke correspondence by
factoring the $7$-division polynomial into all eight cyclic kernels,
including their field norms, and constructing the corresponding
isogenies in Sage. Their target curves and level points agree with
the saved correspondence. Thus this case is impossible.

Finally, Lemma~\ref{lem:decic-isogeny-transfer} gives the cyclic
exclusions of orders $76,80,100$.
\end{proof}

The scripts, finite-field points, class maps, Riemann--Roch witnesses,
and independent replays for this proposition are stored in
\path{continuation_20261029b/}. The accepted computations give
\[
\begin{array}{c|rrrr}
 n&\text{initial patterns}&\text{first filter}&\text{second filter}
       &\text{final patterns}\\\hline
38&233&35&35&0\\
40&781&288&118&0\\
50&137&3&3&0.
\end{array}
\]
The final column uses only completed certificates. Slow alternative
checks that reached their time limits are retained as unsuccessful
probes and are not used in the proof.
